
\documentclass{article}
\usepackage[utf8]{inputenc}
\usepackage[portuguese]{babel}
\usepackage{booktabs}
\usepackage{hyperref}

\title{Análise Linguística: harry.txt}
\author{Gonçalo Magalhães}
\date{\today}

\begin{document}

\maketitle

\begin{abstract}
Frases mais relevantes extraídas através do modelo n-grams:
\begin{itemize}
\item A casa dos Dursley tinha quatro quartos: um para tio Válter e tia Petúnia, um para hóspedes (em geral a irmã de tio Válter, Guida), um onde Duda dormia e um onde Duda guardava todos os brinquedos e pertences que não cabiam no primeiro quarto.
\item Passara a vida dominado por Duda e infernizado pela tia Petúnia e pelo tio Válter; se era realmente um bruxo, por que eles não tinham se transformado em sapos toda vez que tentaram prendê-lo no armário?
\item Enquanto tio Válter dava telefonemas furiosos para o correio e a leiteria tentando encontrar alguém a quem se queixar, tia Petúnia picava as cartas no processador de alimentos. – Mas quem é que quer falar tanto assim com você ?
\end{itemize}
\end{abstract}

\section{Entidades Nomeadas (NER)}
\begin{itemize}
\item \textbf{MISC}: Visite pottermore.com  Para Jessica
\item \textbf{LOC}: Alfeneiros
\item \textbf{MISC}: Harry Potter characters
\item \textbf{MISC}: CAPÍTULO CINCO
\item \textbf{PER}: Continue
\item \textbf{ORG}: TRÊS
\item \textbf{MISC}: Sra. Potter
\item \textbf{PER}: Editora Rocco Ltda
\item \textbf{PER}: Pottermore Limited
\item \textbf{ORG}: DEZESSETE
\item \textbf{LOC}: QUATRO
\item \textbf{MISC}: Harry... 
\item \textbf{MISC}: Harry
\item \textbf{MISC}: Harry Potter
\item \textbf{PER}: Eduardo
\item \textbf{MISC}: Potter
\item \textbf{MISC}: Alfeneiros
\item \textbf{PER}: Sr. Dursley
\item \textbf{MISC}: Le Livros
\item \textbf{ORG}: SETE
\item \textbf{ORG}: DOZE
\item \textbf{ORG}: CATORZE
\item \textbf{LOC}: NOVE
\item \textbf{PER}: Potter
\item \textbf{PER}: Biz
\item \textbf{MISC}: OITO
\item \textbf{PER}: Duda
\item \textbf{MISC}: ONZE
\item \textbf{ORG}: Editora Rocco Ltda
\item \textbf{PER}: J.K. Rowling
\item \textbf{ORG}: DEZESSEIS
\item \textbf{LOC}: Sr.
\item \textbf{MISC}: Direitos Reservados ©
\item \textbf{MISC}: Conteúdo
\item \textbf{ORG}: DOIS
\item \textbf{PER}: Sr. Dursley
\item \textbf{ORG}: Grunnings
\item \textbf{LOC}: Dursley
\item \textbf{PER}: Lia Wyler 
\item \textbf{ORG}: Warner Bros
\item \textbf{MISC}: Título Original: Harry Potter and the Philosopher’s Stone 
\item \textbf{ORG}: TREZE
\item \textbf{LOC}: Desculpe
\item \textbf{ORG}: Ilustrações
\item \textbf{LOC}: Sra. Dursley
\item \textbf{PER}: Direitos Autorais ©
\item \textbf{PER}: Baixelivros.org
\item \textbf{ORG}: Le Livros
\item \textbf{PER}: Anne
\item \textbf{PER}: Dudley
\item \textbf{MISC}: QUINZE
\item \textbf{PER}: Ernesto
\item \textbf{PER}: Mary GrandPré ©
\item \textbf{PER}: Di
\item \textbf{PER}: Dursley
\item \textbf{LOC}: Pestinha
\item \textbf{LOC}: Brasil
\item \textbf{MISC}: DADOS DE COPYRIGHT Sobre
\item \textbf{MISC}: LeLivros
\item \textbf{PER}: Esquecera
\end{itemize}
\clearpage

\section{Análise de Padrões Gramaticais}
\begin{table}[h]
\centering
\begin{tabular}{|l|p{8cm}|}
\hline
\textbf{Tipo de Padrão} & \textbf{Trecho Extraído} \\ \hline
ACAO & Dursley tinha \\ \hline
OBJETO & Duda dormia \\ \hline
ACAO & Duda guardava \\ \hline
MOVIMENTO & transformado em sapos \\ \hline
MOVIMENTO & prendê-lo no armário \\ \hline
ACAO & Válter dava \\ \hline
ACAO & Petúnia picava \\ \hline
OBJETO & picava as cartas \\ \hline

\end{tabular}
\end{table}

\clearpage

\begin{thebibliography}{9}
\bibitem{fonte} Autor Desconhecido, \textit{harry.txt}, Edição Digital.
\end{thebibliography}

\end{document}
